\section{Forelesning 3: \emph{Digitale forretningsmodeller}}
\label{sec:lec03}
\settermlabel{terms:lec03}

\subsection*{Oversikt}
Forelesning 3 introduserer forretningsmodellen som det strategiske bindeleddet
mellom virksomhetens posisjon (\autoref{sec:lec01}) og dens digitale valg.
Hovedkilden er \modelref{weill-woerner}{Weill \& Woerner (2018)}, \emph{What's Your Digital Business Model?},
som gir et todimensjonalt kart over digitale forretningsmodeller. Forelesningen
introduserer også \modelref{bmc}{Business Model Canvas} (Osterwalder \& Pigneur 2010)
som et analytisk -- ikke strategisk -- verktøy, samt \term{Digital Value Chain
(DVC)} fra TEBR.

\subsection{Hovedteori}

\subsubsection*{Hva er en forretningsmodell? (\modelref{cvp}{Johnson et al.})}
Johnson et al.\ (2008) bygger forretningsmodellen rundt fire elementer:
\begin{itemize}
  \item \term{Customer Value Proposition} -- Target Customer + Job to be
    done + Offering: \emph{hvem} løses \emph{hvilket} problem for, og
    \emph{med hva}.
  \item \term{Profit Formula} -- inntektsmodell, kostnadsstruktur og
    marginlogikk.
  \item \term{Key Resources} -- de avgjørende ressursene (data, merkevare,
    teknologi, mennesker).
  \item \term{Key Processes} -- de gjentakbare prosessene som leverer CVP-en.
\end{itemize}
Forretningsmodellen beskriver hvordan virksomheten skaper, leverer og fanger
verdi. Den er fundamentet som strategi bygger på -- ikke et substitutt for
strategi.

\begin{figure}[H]
\centering
\includegraphics[width=0.85\linewidth]{lec03-johnson-four-building-blocks.png}
\caption{Johnson et al.\ (2008): forretningsmodellens fire byggeklosser med
CVP i sentrum.}
\end{figure}

\begin{definition}{Digital forretningsmodell}
En forretningsmodell der digital teknologi er sentral for verdiskaping,
verdilevering og/eller verdifangst. Karakteriseres av mer åpne nettverk,
sterkere kunderelasjoner og dypere kundekunnskap enn klassiske modeller.
\end{definition}

\subsubsection*{Justering vs.\ reinvention (Johnson)}
Ikke alle modellendringer er like dyptgripende. Johnson skiller mellom
\emph{justeret forretningsmodel} (inkrementell tilpasning) og \emph{ny
forretningsmodel} (gjennomgripende reinvention). Fire diagnostiske spørsmål
hjelper avgjøre når reinvention er berettiget.

\begin{figure}[H]
\centering
\includegraphics[width=0.85\linewidth]{lec03-adjustment-vs-reinvention-curve.png}
\caption{Johnsons kurve: \emph{adjustment} (inkrementell) vs.\
\emph{reinvention} (ny modell) over tid.}
\end{figure}

\subsubsection*{\modelref{weill-woerner}{Weill \& Woerner (2018)}: To dimensjoner}
Rammeverket plasserer virksomheter langs to akser:
\begin{enumerate}
  \item \term{Knowledge of end customer} -- hvor godt kjenner virksomheten
    sluttkunden? (Partial $\rightarrow$ Complete)
  \item \term{Business design} -- hvordan er verdinettverket organisert?
    (Value chain $\rightarrow$ Ecosystem)
\end{enumerate}

Dette gir fire grunnmodeller:

\begin{definition}{De fire digitale forretningsmodellene}
\begin{enumerate}
  \item \term{Supplier} -- selger gjennom andre, lav direkte kundekunnskap,
    operer i kontrollert verdikjede. Eks.: tradisjonell merkevareleverandør.
  \item \term{Omnichannel business} -- eier kunderelasjonen på tvers av
    fysiske og digitale kanaler, moderat-til-høy kundekunnskap, kontrollert
    kjede. Eks.: Coop, banker, store retailere.
  \item \term{Modular producer} -- plug-and-play-kapabiliteter solgt til
    andres økosystemer. Eks.: PayPal, Stripe; potensielt \emph{Lobyco}.
  \item \term{Ecosystem driver} -- hub i et bredt økosystem, høy
    kundekunnskap, koordinerer flere aktører. Eks.: Amazon, Alibaba, Apple.
\end{enumerate}
\end{definition}

\begin{figure}[H]
\centering
\includegraphics[width=0.9\linewidth]{lec03-weill-woerner-digital-business-model-framework.png}
\caption{Weill \& Woerners (2018) Digital Business Model Framework: 2$\times$2
av \emph{Business design} $\times$ \emph{Knowledge of end customer} med de
fire arketypene Supplier, Modular producer, Omnichannel, Ecosystem driver.}
\end{figure}

\subsubsection*{Tre dimensjoner for konkurransefortrinn}
Weill \& Woerner identifiserer tre dimensjoner virksomheter må mestre:
\begin{itemize}
  \item \term{Content} -- hva tilbys digitalt (produkter, informasjon, tjenester).
  \item \term{Customer experience} -- hvordan digital interaksjon oppleves.
  \item \term{Platform} -- felles digital infrastruktur på tvers av aktiviteter.
\end{itemize}

\subsubsection*{\modelref{bmc}{Business Model Canvas} (Osterwalder \& Pigneur 2010)}
\modelref{bmc}{BMC} kartlegger forretningsmodellen i 9 byggeklosser. \keypoint{Value
Proposition} står i sentrum og kobler kundesiden (segmenter, kanaler,
relasjoner, inntektsstrømmer) med infrastruktursiden (nøkkelpartnere,
nøkkelaktiviteter, nøkkelressurser, kostnadsstruktur).

\begin{figure}[H]
\centering
\includegraphics[width=0.9\linewidth]{lec03-business-model-canvas.png}
\caption{Business Model Canvas -- 9 byggeklosser med Value Proposition i sentrum.}
\end{figure}

\keypoint{BMC er et analytisk verktøy, ikke en strategi i seg selv.}
Den beskriver en forretningsmodell i et snitt -- ikke konkurransedynamikk,
ikke endring over tid, ikke avveininger.

\begin{figure}[H]
\centering
\includegraphics[width=0.9\linewidth]{lec03-business-model-canvas-airbnb.png}
\caption{Eksempel: BMC anvendt på Airbnb. Brukes som referanse når man
fyller ut canvas for en digital forretningsmodell.}
\end{figure}

\subsubsection*{DVC Framework (TEBR/Rosenstand \& Baiyere)}
DVC er et komplementært rammeverk som diagnostiserer \emph{hvor} digital verdi
kan oppstå:
\begin{itemize}
  \item \term{Experiences} -- digitale interaksjoner som forbedrer kundeopplevelsen.
  \item \term{Relationships} -- digitale kanaler som utdyper kunderelasjoner.
  \item \term{Relevance} -- digital verdi tilpasset kundens situasjon.
  \item \term{Digital Infrastructure} -- back-end-kapabiliteter.
  \item \term{Digital Outputs} -- målbare digitale resultater.
\end{itemize}
\keypoint{Kritisk poeng: Digital outputs $\neq$ digital outcomes. DVC alene
kan gjøre digital aktivitet til å fremstå verdifull bare fordi den
\emph{eksisterer}. Kombiner med DBM for å sjekke om verdien faktisk
\emph{fanges}.}

\subsection{Sentrale fagbegreper}\label{terms:lec03}
\begin{itemize}
  \item \term{Value proposition} -- løftet til kunden som forklarer hvorfor de
    velger oss fremfor alternativer.
  \item \term{Value creation} -- hvem som får verdi og hvordan.
  \item \term{Value capture} -- hvordan virksomheten fanger økonomisk verdi.
  \item \term{Customer knowledge} -- dybden i forståelsen av sluttkunden.
  \item \term{Ecosystem} vs.\ \term{value chain} -- nettverksbasert vs.\ lineær.
  \item \term{Adjustment} vs.\ \term{reinvention} -- inkrementell vs.\
    fundamental modellendring.
\end{itemize}

\subsection{Modeller -- bruksveiledning}

\begin{enumerate}
  \item Bruk \emph{Johnson} til å definere hva en forretningsmodell består av
    før du analyserer den.
  \item Bruk \emph{Weill \& Woerner} til å klassifisere virksomheten (en av
    fire) og argumentere for strategiske grep \emph{mellom} kvadrantene.
  \item Bruk \emph{BMC} som strukturen i en konkret analyse eller en
    fremtidig anbefalt modell.
  \item Bruk \emph{DVC} til å finne \emph{hvor} digital verdi kan oppstå --
    deretter kontroller med DBM at den \emph{fanges} i modellen.
  \item Bruk \emph{adjustment/reinvention}-spørsmålene for å forsvare om en
    case krever liten eller stor endring.
\end{enumerate}

\subsection{Eksempler og caser}
\begin{eksempelboks}{Coop som omnichannel}
Klassisk \emph{omnichannel business}. Etter lukkingen av Coop.dk MAD (2023)
og Coop.dk-webshop (januar 2025) er posisjonen entydig store-first
omnichannel. Customer knowledge er moderat-til-høy (medlemsdata fra Coop
Medlem). Business design er value-chain-orientert, ikke ecosystem.
\end{eksempelboks}

\begin{eksempelboks}{Lobyco som modular producer}
Potensiell \emph{modular producer} hvis lojalitets- og betalings-
infrastrukturen selges til andre retailere. Lobyco gjør Coops digitale
kapabilitet til en separat plug-and-play-tjeneste -- et eksempel på
strategisk bevegelse mellom kvadrantene i Weill \& Woerner.
\end{eksempelboks}

\begin{eksempelboks}{Amazon som ecosystem driver}
Prototypen på \emph{ecosystem driver}: høy kundekunnskap + bred koordinering
av tredjeparter (Marketplace, AWS, Prime, Alexa). Eksempel som tradisjonelle
retailere (Coop) blir målt mot, men ikke nødvendigvis bør imitere.
\end{eksempelboks}

\subsection{Eksamensrelevans}
\begin{itemize}
  \item Kunne navngi og forklare de fire DBM-typene med eksempler.
  \item Kunne plassere en konkret virksomhet (f.eks.\ Coop) i matrisen og
    forsvare valget med evidens.
  \item Forklare hvorfor BMC ikke er en strategi.
  \item Vise hvordan DVC + DBM brukes sammen (DVC: \emph{hvor} kan verdi
    skapes; DBM: blir verdien \emph{fanget} på en måte som passer modellen?).
  \item Bruke Johnsons adjustment/reinvention-spørsmål til å forsvare
    omfanget av en anbefalt endring.
  \item Knytte til \autoref{sec:lec04} (digital transformasjon) og
    \autoref{sec:lec06} (økosystemer).
\end{itemize}
